
\section{基变换下保持的性质}

态射的一切「合理」性质都在\trBaseChange{}下被保持（参看 \S8.1(ii)）。我们将讨论这一点，并在 \S10.5.1 中解释如何修补那些不合此模式的性质。

我们已经证明过：「\trOpenEmbedding」这一观念在\trBaseChange{}下被保持（\cref{ex:open-embeddings-base-change}）。我们的做法是显式地描述出一个\trOpenEmbedding{}的\trFiberProduct{}是什么：若 $Y\hookrightarrow Z$ 是一个\trOpenEmbedding，而 $\psi\colon X\to Z$ 是任意\trMorphism，则我们验证过：$X$ 的开子概形 $\psi^{-1}(Y)$ 满足\trFiberProduct{}的\trUniversalProperty。

我们还证明过：「\trClosedEmbedding」这一观念在\trBaseChange{}下被保持（\S10.2(3)）。换言之，给定一个\trCartesianDiagram
\[
\begin{tikzcd}[row sep=1.9em, column sep=3.2em]
W \ar[r] \ar[d] & Y \ar[d, "\text{闭嵌入}"] \\
X \ar[r] & Z
\end{tikzcd}
\]
其中 $Y\hookrightarrow Z$ 是\trClosedEmbedding，则 $W\to X$ 也是。

\textbf{10.4.A.（容易的练习）}\quad 证明：\trClosedEmbedding{}与\trLocally\trClosedEmbedding{}在 \S8.1 的意义下都是「合理」的态射类。

\textbf{10.4.B.（练习）}\quad 设 $X\to Z$ 与 $Y\to Z$ 都是\trLocally\trClosedEmbedding。(a)~证明：$X\times_Z Y\to Z$ 是一个\trLocally\trClosedEmbedding。(b)~由此我们定义 $Z$ 的任意两个\trLocally\trClosed\trSubscheme{}的\textbf{（概形论）交}（(scheme-theoretic) intersection）。请验证：这与先前对 $Z$ 的两个\trOpenSubscheme{}之交的定义、以及先前对 $Z$ 的两个\trClosedSubscheme{}之交的定义相一致。

\textbf{10.4.C.（容易的练习）}\quad 证明：\trLocallyPrincipal{}闭子概形（定义~9.5.1）经\trPullback{}仍是\trLocallyPrincipal{}闭子概形。

类似地，另一些重要性质也在\trBaseChange{}下被保持。

\textbf{10.4.D.（练习）}\quad 证明：下列\trMorphism{}性质在\trBaseChange{}下被保持。
\begin{enumerate}[label=(\alph*)]
\item \trQuasicompact
\item \trAffine\trMorphism
\item \trFinite
\item \trIntegral
\item \trLocally\trFiniteType
\item \trFiniteType
\item[\hard~(g)] \trLocally\trFinitelyPresented
\end{enumerate}

\textbf{10.4.E.（练习）}\quad 证明：上一道练习中的那些\trMorphism{}性质，在 \S8.1 的意义下都是「合理」的。

\vkexsync{F}
\begin{exercise}\label{ex:quasifinite-preserved-base-change}
\begin{enumerate}[label=(\alph*)]
  \item[\hard~(a)] 证明：「\trQuasifinite\trMorphism」（有限型 $+$ 有限的\trFiber，定义~8.3.11）这一观念在\trBaseChange{}下被保持。（警告：「有限的\trFiber」这一观念在\trBaseChange{}下并\emph{不}被保持。$\Spec\bar{\mathbb{Q}}\to\Spec\mathbb{Q}$ 有有限的\trFiber，但 $\Spec\bar{\mathbb{Q}}\otimes_{\mathbb{Q}}\bar{\mathbb{Q}}\to\Spec\bar{\mathbb{Q}}$ 对 $\Gal(\bar{\mathbb{Q}}/\mathbb{Q})$ 的每个元素恰有一个点；见\mbox{练习~10.2.F}。）提示：化归到 $\Spec\scr{A}\to\Spec B$ 的情形。再化归到 $\phi\colon\Spec\scr{A}\to\Spec k$ 的情形。由\cref{ex:quasifinite-to-field-finite}，这样的 $\phi$ 实际上是有限的，而有限性在\trBaseChange{}下被保持。
  \item 证明：「\trQuasifinite\trMorphism」这一观念在 \S8.1 的意义下是「合理」的。
\end{enumerate}
\end{exercise}

\textbf{10.4.G.（练习）}\quad (a)~证明：满性在\trBaseChange{}下被保持。（\textbf{满性}（surjectivity）取其通常的含义：作为集合的映射是满的。）你最终也许会去证明：对任意包含 $k_3$ 的域 $k_1$ 与 $k_2$，都有 $k_1\otimes_{k_3}k_2$ 非零，并借\trZornsLemma{}找出 $k_1\otimes_{k_3}k_2$ 中的一个\trMaximalIdeal。(b)~证明：满态射构成一个「合理」的类（\S8.1）。

\textbf{10.4.1.}\quad 另一方面，单性在\trBaseChange{}下不被保持——双射 $\Spec\mathbb{C}\to\Spec\mathbb{R}$ 便是见证：它经 $\Spec\mathbb{C}\to\Spec\mathbb{R}$ 作\trBaseChange{}后便失去单性（见例~10.2.3）。这件事可以补救（见 \S10.5.1）。

\textbf{10.4.H.（练习：参看\mbox{练习~10.2.E}）}\quad 设 $X$ 与 $Y$ 都是整的\trFiniteType{} $\bar{k}$-概形。证明：$X\times_{\bar{k}}Y$ 是一个整的\trFiniteType{} $\bar{k}$-概形。（一旦我们定义了「\trVariety」，这将成为那条重要事实：\trAlgClosedField{}上\trIrreducible{}簇的积仍是\trIrreducible{}簇；参看命题~11.2.11。\trBaseField{} $\bar{k}$ 是\trAlgClosed{}的这一点是要紧的；见 \S10.5.1。改进的版本见\mbox{练习~10.5.O}与\mbox{练习~10.5.R}。）提示：化归到 $X$ 与 $Y$ 都是仿射的情形，比如说 $X=\Spec\scr{A}$、$Y=\Spec B$，其中 $\scr{A}$ 与 $B$ 都是\trIntegralDomain。你不妨先翻到\mbox{容易的练习~10.5.N}看看怎么做。设在 $\scr{A}\otimes_{\bar{k}}B$ 中 $\bigl(\sum a_i\otimes b_i\bigr)\bigl(\sum a'_j\otimes b'_j\bigr)=0$，其中 $a_i$、$a'_j\in\scr{A}$，$b_i$、$b'_j\in B$，而 $\{b_i\}$ 与 $\{b'_j\}$ 都在 $\bar{k}$ 上线性无关，且 $a_1$ 与 $a'_1$ 非零。证明 $D(a_1a'_1)\subset\Spec\scr{A}$ 非空。由弱\trNullstellensatz~3.2.5，在 $D(a_1a'_1)$ 中存在\trMaximalIdeal{} $\mathfrak{m}\subset\scr{A}$ 使得 $\scr{A}/\mathfrak{m}=\bar{k}$。对 $\mathfrak{m}$ 取剩余，推出在 $B$ 中 $\bigl(\sum\overline{a}_i\otimes b_i\bigr)\bigl(\sum\overline{a}'_j\otimes b'_j\bigr)=0$，其中横线表示按 $\mathfrak{m}$ 取剩余。证明这与 $B$ 是\trIntegralDomain{}这一事实相矛盾。
